\subsection{Microscopic parent constraints and native seam boundaries}
\label{sec:peps_native_parent_seam_cuts}

The parent-Hamiltonian argument in \cite{Schuch2010PEPS} uses the regional
intersection and closure statements of Theorems~5.4 and~5.5. The auxiliary
results below pass directly from microscopic plaquette constraints to
boundary tensors on the original lattice.
% Scope: native simple tori, both periods at least three. No identification
% with the separate four-block closure space is claimed in this subsection.

\begin{definition}[Regional cuts in the original graph]%
    \label{def:peps_native_regional_cut}
    \lean{TNLean.PEPS.graphEdgeTail,TNLean.PEPS.graphEdgeHead,
        TNLean.PEPS.graphBondAlphabet,TNLean.PEPS.graphIncidentEndpoint,
        TNLean.PEPS.graphDependentTensor,TNLean.PEPS.regionExteriorCut,
        TNLean.PEPS.regionCutBoundaryLabel,TNLean.PEPS.graphDependentTensor_diagonal,
        TNLean.PEPS.network_graphDependentTensor_one,TNLean.PEPS.mem_regionExteriorCut_of_boundary,
        TNLean.PEPS.mem_regionExteriorCut_of_outside,TNLean.PEPS.regionCutBoundaryLabel_diagonal,
        TNLean.PEPS.graphEndpointIncident,TNLean.PEPS.graphIncidentEndpoint_endpointIncident}
    \leanok
    Let $\Gamma=(V,E)$ be a finite simple graph with ordered vertices, where
    each $e\in E$ has tail $t(e)<h(e)$ and alphabet $X_e=\{0,\ldots,D_e-1\}$,
    and each site has tensor $A_v$ with physical codomain $\C^d$ (until
    Lemma~\ref{thm:peps_nonuniform_parent_seam_boundary}). For $R\subseteq V$,
    define $C_R=\{e\in E:\{t(e),h(e)\}\not\subseteq R\}=E\setminus E_R$.
    Identifying each incident edge at $v$ with its endpoint label matches the
    original tensor coordinates with independent endpoint coordinates, so
    contracting identity bonds yields $\Psi(A)$ and diagonal endpoint
    assignments restrict to the original crossing labels on $\partial E_R$.
    Write $\mathcal S_{A,C}$ for the range of the cut contraction against a
    joint boundary tensor on the doubled incidences of $C$.
\end{definition}

\begin{lemma}[Parent constraints yield regional cut boundaries]%
    \label{thm:peps_native_parent_regional_cut}
    \lean{TNLean.PEPS.regionParentGroundSpace_le_iInf_cutSpace,
        TNLean.PEPS.exists_cutBoundary_of_mem_regionParentKernel,
        TNLean.Lattice.DependentBondNetwork.headTailConfigEquiv,
        TNLean.Lattice.DependentBondNetwork.diagonalEndpointConfig,
        TNLean.Lattice.DependentBondNetwork.sum_identityWeight_mul,
        TNLean.Lattice.DependentBondNetwork.cutIdentityWeight,
        TNLean.Lattice.DependentBondNetwork.cutIdentityWeight_mul_interior,
        TNLean.Lattice.DependentBondNetwork.sum_cutIdentityWeight_mul,
        TNLean.PEPS.sum_regionBoundaryLabel_mul_eq_openRegionMap,
        TNLean.PEPS.outsidePhysicalToComplement,TNLean.PEPS.graphRegionalCutBoundary,
        TNLean.PEPS.liftedCutMap_graphRegionalCutBoundary,
        TNLean.PEPS.mem_liftedCutSpace_of_region_slices,
        TNLean.PEPS.regionParentGroundSpace_le_iInf_liftedCutSpace}
    \leanok
    \uses{def:peps_native_regional_cut}
    Suppose $D_e>0$ for every $e\in E$ and the finite region family
    $(R_i)_i$ covers $V$. For $H_i\geq0$ with $\ker H_i=\mathcal G_{R_i}(A)$
    and $H=\sum_i H_i\otimes I_{V\setminus R_i}$, one has
    $\ker H\subseteq\bigcap_i\mathcal S_{A,C_{R_i}}$, without requiring
    injectivity or uniform bond dimensions.
\end{lemma}
\begin{proof}\leanok
    Since $H_i\ge 0$, every $\psi\in\ker H$ has regional slices in
    $\mathcal G_{R_i}(A)$ with virtual boundaries indexed by the exterior
    physical configuration on $V\setminus R_i$. Closing each exterior bond
    $e\in E_{R_i^c}\subseteq C_{R_i}$ with $\mathrm{id}_{X_e}$ multiplies the
    original contraction by $\prod_{e\in E_{R_i^c}}D_e>0$; dividing by this
    factor places $\psi$ in the lifted regional cut space. Since $(R_i)_i$
    covers $V$, every single-site slice lies in $\operatorname{im}A_v$, so
    local linear retractions on $V\setminus R_i$ absorb the exterior physical
    indices into the joint virtual boundary on $C_{R_i}$, giving
    $\psi\in\mathcal S_{A,C_{R_i}}$.
\end{proof}

\begin{lemma}[Coherent coefficients from the physical parent]%
    \label{thm:peps_native_parent_coefficients}
    \lean{TNLean.PEPS.exists_regionParent_commonInverse,
        TNLean.PEPS.exists_regionParent_bondCoefficients,
        TNLean.PEPS.bondDim_ne_zero_of_isSemiRegular}
    \leanok
    \uses{thm:peps_native_parent_regional_cut}
    Let $G$ be finite, let each $e\in E$ carry a semi-regular representation
    $U_e$, and let each $A_v$ be $G$-injective for its incident action. If the
    nonempty finite family $(R_i)_i$ covers $V$ and $E$ ($E=\bigcup_i E_{R_i}$),
    every $\psi\in\mathcal P_{\mathcal R}(A)$ admits canonical coordinates
    $\xi$ with $\psi=(\bigotimes_v A_v)\xi$ and
    \begin{align}
        \xi          & =\sum_{p:E\to G}c_p(\xi)b_p,\nonumber \\
        c_p(\xi)\ne0 & \ \Longrightarrow
        \forall i\ \exists q_i:V\to G\ \forall e\notin C_{R_i},\quad
        p_e=q_i(h(e))q_i(t(e))^{-1},
        \label{eq:peps_native_parent_coefficients}
    \end{align}
    where $b_p=\bigotimes_{e\in E}\operatorname{vec}(U_e(p_e))$ and $c_p$ is
    the product trace-dual functional, obtained via a single family of local
    inverses.
\end{lemma}
\begin{proof}\leanok
    Semi-regularity forces $D_e>0$ for all $e\in E$. Applying one family of
    local $G$-injective inverses to the cut boundaries from
    Lemma~\ref{thm:peps_native_parent_regional_cut} produces $\xi$ with
    $\psi=(\bigotimes_v A_v)\xi$ and places every internal bond $e\in E_{R_i}$
    in $\spn\{\operatorname{vec}(U_e(g)):g\in G\}$. Edge coverage and the
    product-range criterion give the expansion $\xi=\sum_p c_p(\xi)b_p$, where
    the trace-dual functionals $c_p$ force $p_e=q_i(h(e))q_i(t(e))^{-1}$ on
    $E_{R_i}=E\setminus C_{R_i}$ whenever $c_p(\xi)\ne 0$.
\end{proof}

\begin{definition}[Native graph seams]%
    \label{def:peps_native_graph_seams}
    \lean{TNLean.PEPS.torusGraphSeamCut,
        TNLean.PEPS.torusRightEdge_mem_graphSeamCut,
        TNLean.PEPS.torusUpEdge_mem_graphSeamCut}
    \leanok
    On $V=\Z_w\times\Z_h$ with $w,h\geq3$, the seam cut $C_{c,r}$ consists of
    all horizontal bonds $(x,y)\to(x+1,y)$ with $x+1=c$ and all vertical bonds
    $(x,y)\to(x,y+1)$ with $y+1=r$.
\end{definition}

\begin{lemma}[Microscopic plaquette parents have every native seam boundary]%
    \label{thm:peps_native_parent_seam_boundary}
    \lean{TNLean.PEPS.torusPlaquetteParentGroundSpace_le_graphSeamCut,
        TNLean.PEPS.exists_graphSeamBoundary_of_mem_torusPlaquetteParentKernel,
        TNLean.PEPS.regularWalkHolonomy_eq_relative_of_region,
        TNLean.PEPS.regularWalkHolonomy_eq_one_of_region,TNLean.PEPS.IsTorusFlat.translate,
        TNLean.PEPS.exists_vertexGauge_off_torusGraphSeamCut,
        TNLean.PEPS.averagingProjector_bondProduct_mem_graphSeamCut}
    \leanok
    \uses{thm:peps_native_parent_coefficients,def:peps_native_graph_seams}
    On $V=\Z_w\times\Z_h$ ($w,h\ge 3$) with semi-regular $U_e$ and
    independently $G$-injective $A_v$, any positive microscopic plaquette
    parent $H$ with exact open-plaquette kernels satisfies
    $\ker H\subseteq\mathcal S_{A,C_{c,r}}$ for every $(c,r)\in V$.
\end{lemma}
\begin{proof}\leanok
    Elementary plaquettes cover $V$ and $E$. For each $p$ with $c_p(\xi)\ne 0$
    in~(\ref{eq:peps_native_parent_coefficients}), the local coboundary relations
    $p_e=q_i(h(e))q_i(t(e))^{-1}$ telescope around each plaquette to give
    trivial plaquette holonomy. Translating the spanning-tree gauge to origin
    $(c,r)$ yields $g:V\to G$ with $g_{h(e)}p_eg_{t(e)}^{-1}=1$ for all
    $e\notin C_{c,r}$. Applying the product of local group averages converts
    each $b_p$ (and hence $\xi$) into a $C_{c,r}$-cut contraction, and
    $\bigotimes_v A_v$ maps it to $\psi\in\mathcal S_{A,C_{c,r}}$.
\end{proof}

\begin{lemma}[Independent physical dimensions and the entire ambient kernel]%
    \label{thm:peps_nonuniform_parent_seam_boundary}
    \lean{TNLean.PEPS.exists_nonuniform_seamBoundaries_of_mem_torusPlaquetteParentKernel,
        TNLean.PEPS.physicalPaddingEmbedding,TNLean.PEPS.graphPhysicalPaddingTensor,
        TNLean.PEPS.physicalRestrictionProduct_padding,
        TNLean.PEPS.padding_restriction_eq_self_of_mem_regionParentGroundSpace,
        TNLean.PEPS.nonuniform_physicalRepresentative_unique,
        TNLean.PEPS.physicalPaddingDimension,TNLean.PEPS.physicalPaddingMatrix,
        TNLean.PEPS.physicalRestrictionMatrix,TNLean.PEPS.physicalRestrictionMatrix_mul_padding,
        TNLean.PEPS.physicalPaddingMatrix_injective,TNLean.PEPS.physicalPaddingProduct_injective,
        TNLean.PEPS.graphDependentTensor_physicalPadding,
        TNLean.PEPS.physicalMapSite_restrict_graphPhysicalPaddingTensor,
        TNLean.PEPS.isGInjective_graphPhysicalPaddingTensor,
        TNLean.PEPS.cutSpace_graphPhysicalPaddingTensor,
        TNLean.PEPS.restriction_cutMap_graphPhysicalPaddingTensor}
    \leanok
    \uses{thm:peps_native_parent_seam_boundary}
    For site-dependent finite physical alphabets $Y_v$, let
    $E_v:\C^{Y_v}\to\C^{\sum_w|Y_w|}$ be coordinate zero-extension with left
    inverse $R_vE_v=I$, and form $B_v=E_vA_v$ with positive plaquette parent
    $H_B$. Every $\Psi\in\ker H_B$ (without prior restriction to
    $\bigotimes_v\operatorname{im}E_v$) has a unique $\psi\in\bigotimes_v\C^{Y_v}$
    satisfying
    \begin{align}
        \Psi & =(\bigotimes_v E_v)\psi,
             & \psi                                    & =(\bigotimes_v R_v)\Psi,\nonumber \\
        \psi & \in\bigcap_{c,r}\mathcal S_{A,C_{c,r}}.
        \label{eq:peps_nonuniform_parent_recovery}
    \end{align}
\end{lemma}
\begin{proof}\leanok
    Since $R_vE_v=I$, each $E_v$ is injective, $B_v=E_vA_v$ remains
    $G$-injective, and cut contractions intertwine with $\bigotimes_v E_v$ and
    $\bigotimes_v R_v$. By Lemma~\ref{thm:peps_native_parent_seam_boundary}
    applied to $B$, every $\Psi\in\ker H_B$ lies in
    $\bigcap_{c,r}\mathcal S_{B,C_{c,r}}=(\bigotimes_v E_v)(\bigcap_{c,r}\mathcal S_{A,C_{c,r}})$,
    and $(\bigotimes_v R_v)(\bigotimes_v E_v)=I$ gives uniqueness and~(\ref{eq:peps_nonuniform_parent_recovery}).
\end{proof}
